\documentclass[UTF8,11pt,fontset=fandol]{ctexart}
\usepackage[a4paper,margin=21mm,top=20mm,bottom=20mm]{geometry}
\usepackage{xcolor,tabularx,array,enumitem,fancyhdr,hyperref,tikz}
\usetikzlibrary{arrows.meta,positioning}
\definecolor{green}{HTML}{087B62}
\definecolor{ink}{HTML}{182C37}
\definecolor{muted}{HTML}{5D6E78}
\definecolor{pale}{HTML}{EDF5F2}
\definecolor{line}{HTML}{D8E5DF}
\hypersetup{colorlinks=true,urlcolor=green,linkcolor=green,pdftitle={楼里软硬件产品设计与商业化讨论稿},pdfauthor={楼里项目}}
\setlength{\parindent}{0pt}
\setlength{\parskip}{6pt}
\linespread{1.22}
\setlist[itemize]{leftmargin=1.5em,itemsep=3pt,topsep=3pt}
\setlist[enumerate]{leftmargin=1.7em,itemsep=3pt,topsep=3pt}
\renewcommand{\arraystretch}{1.4}
\newcolumntype{Y}{>{\raggedright\arraybackslash}X}
\pagestyle{fancy}\fancyhf{}
\fancyhead[L]{\small\color{muted} 楼里\quad 产品设计与商业化讨论稿}
\fancyhead[R]{\small\color{muted}2026 年 10 月}
\fancyfoot[L]{\footnotesize\color{muted}产品概念与仿真演示阶段}
\fancyfoot[R]{\small\color{green}\thepage}
\renewcommand{\headrulewidth}{0pt}\setlength{\headheight}{15pt}
\ctexset{section={format=\Large\bfseries\color{ink},beforeskip=13pt,afterskip=8pt},subsection={format=\normalsize\bfseries\color{green},beforeskip=10pt,afterskip=4pt}}
\newcommand{\tagline}[1]{{\small\color{green}\bfseries #1}\par}
\newcommand{\lead}[1]{\vspace{4pt}\colorbox{pale}{\parbox{\dimexpr\linewidth-16pt\relax}{\vspace{6pt}\color{ink}#1\vspace{6pt}}}\par\vspace{4pt}}
\newcommand{\page}[2]{\clearpage\tagline{#1}\section*{#2}}
\newcommand{\tbl}[1]{\par\begingroup\small\noindent\begin{tabularx}{\linewidth}{#1}}
\begin{document}
\thispagestyle{empty}
\tagline{LOULI\quad 供团队讨论的产品设计稿}
{\fontsize{32}{40}\selectfont\bfseries\color{ink}楼里医院就诊导航\par}
\vspace{7pt}
{\Large\color{muted}软硬件概念\quad 展示方案\quad 商业化路径\par}
\vspace{12pt}
\lead{\textbf{建议团队统一的方向}\quad 先做能拿出去讲、演示和谈合作的产品概念，用小范围试点验证客户价值。近期把时间投入产品表达、展示故事与合作材料，再依据反馈推进技术实现。}
\subsection*{我们为谁解决什么问题}
首次来院的患者和陪诊家属，需要知道“下一步去哪、怎么走、到哪里等”。医院门诊管理团队需要改善导诊体验；信息科与信息化合作商关注系统接入、部署维护和交付边界。
\subsection*{产品由三部分组成}
\begin{itemize}
\item \textbf{患者端}\quad 科室查询、路线指引、就诊任务清单与返回提醒。
\item \textbf{场馆端}\quad 室内地图、科室与通道信息、路线维护和试点反馈。
\item \textbf{楼里信标}\quad 场馆内的蓝牙定位参考点，配合手机与地图辅助确认位置。
\end{itemize}
\subsection*{现在能展示什么}
仓库已有商场与医院网页仿真，可展示路线规划、定位方法对比、模拟排队与多点就诊流程。信号、轨迹和队列由仿真生成；硬件外观、管理端和真实系统接入属于本稿中的产品设计。
\subsection*{后续洽谈希望得到什么}
一次业务访谈、一次楼层勘察，以及一位试点对接人。双方确认范围后，再决定是否开展短期试点；首轮可用人工配置流程和固定起点演示，之后评估信标验证与系统接入。
\subsection*{拟议商业模式}
先以范围明确的场景调研或试点项目验证价值，再形成部署服务与年度维护方案。与医疗信息化集成商合作时，楼里可作为其患者服务模块。具体价格按场地、功能、接入工作量和验收要求商议。
\vfill
{\small\color{muted}2026 年 10 月 7 日\quad 团队讨论版本 1.0\par
先供队友讨论方向与分工。文中范围、排期和商业路径为建议；后续对外版本依据团队决定调整。}

\page{00\quad 团队决策}{这轮做到什么程度就可以出去谈}
\lead{\textbf{近期完成标准}\quad 队友能统一介绍产品，观众能看懂完整故事，合作方能知道我们希望他提供什么。用这三个标准安排工作。}
\subsection*{三个推进方式与建议}
\tbl{p{28mm}YY}
\textbf{方式}&\textbf{适合做什么}&\textbf{本轮取舍}\\\hline
概念与销售展示&设计文稿、硬件概念、软件演示和洽谈&建议优先，支撑 IDP 汇报与客户访谈\\
真实硬件体验&现成信标、手机扫描、现场验证&作为第二步，取决于队员与设备条件\\
完整系统研发&定制硬件、定位软件、业务系统集成&放在范围与合作资源明确以后\\
\end{tabularx}\endgroup\par
\subsection*{第一轮展示包}
一份产品设计稿、一套两分钟可重复演示的患者故事、一张硬件外观与布设说明，以及一页试点合作摘要。软件复用当前仿真；硬件先用概念图或外观模型解释，能取得现成信标时再增加广播演示。该组合是团队拟做范围，尚未全部完成。
\subsection*{建议分工}
\tbl{p{29mm}Y}
\textbf{商业与统筹}&约访、问题清单、现场讲述、会后反馈与合作推进\\\hline
\textbf{产品与展示}&患者故事、画面文案、演示节奏、截图与录屏\\\hline
\textbf{软件与验证}&复用网页 Demo、故事重置、异常处理与演示稳定性\\\hline
\textbf{硬件与合作}&外观概念、样品询价、配置工具与供应商沟通\\
\end{tabularx}\endgroup\par
人数较少时可以合并角色。具体人名与时间由队友讨论后确定，不在本稿中预设承诺。
\subsection*{建议的团队会议议程}
\begin{enumerate}
\item 是否先以医院流程导航为主故事，商场作为补充展示？
\item 本轮只做概念外观，还是有资源加入真实信标扫描？
\item 谁负责展示与约访，谁负责软件和硬件支持？
\item 最近一次展示是什么时间，先接触哪类合作对象？
\end{enumerate}
\subsection*{讨论后的产出}
确定一个故事、一份展示范围、每项工作负责人和一次演示检查时间。技术验证另列任务；首轮客户沟通收集的需求再进入下一轮开发。

\page{01\quad 商业定位}{先卖清楚一段就诊流程}
\lead{\textbf{核心判断}\quad 首轮销售的对象，是一个可评估的患者服务方案及其试点合作。技术路线服务于这个方案。}
\subsection*{为什么从医院切入}
医院就诊常涉及多个地点与先后步骤，路线需要结合具体业务流程。已有院方公开招标文件包含目的地导航、信息提醒、预计步行时间等要求，说明这一类需求曾进入实际采购范围。[2] 该历史案例不能证明某家医院现在有预算，也不能替代本地访谈。
\subsection*{买方和使用者要分开理解}
\tbl{p{25mm}YY}
\textbf{角色}&\textbf{最关心的问题}&\textbf{面谈希望确认的事}\\\hline
患者与家属&能不能找到下一站，少走冤枉路&当前怎样找路，在哪些节点最困惑\\
门诊办与导诊团队&能否改善服务，减少反复问路&最值得试点的流程，业务负责人是谁\\
信息科&能否接入，谁维护，怎样验收&现有系统与地图，授权和采购路径\\
集成商&是否补充其方案，交付是否可控&模块范围、渠道分工、售后责任\\
\end{tabularx}\endgroup\par
\subsection*{三个可选产品方向}
\begin{itemize}
\item \textbf{医院流程导航，建议优先}\quad 将检查、候诊、复诊和取药串联。价值需要业务数据支撑，系统对接较重；适合与院方或集成商共创。
\item \textbf{商场与停车场导航}\quad 演示直观、场景容易理解。付费理由应落到具体服务问题，例如反向寻车；当前保留为扩展场景。
\item \textbf{信标硬件销售}\quad 可借助供应链较快提供样机，但需要验证采购需求、维护成本与产品差异。首轮更适合作为方案的一部分。
\end{itemize}
\subsection*{首个场景的选择标准}
选一层楼内、经过两至四个服务点、业务负责人可参与的流程。优先包含真实的找路困难，且允许现场观察；不要求首轮读取患者病历或接入生产系统。科室顺序和检查前置条件由院方确认。
\subsection*{首轮商业验证问题}
是否存在明确痛点？谁愿意组织试点？改善到什么程度值得继续？从哪项预算支付？谁拥有后续决策权？这五个答案，比早期三年收入预测更能指导下一步。

\page{02\quad 硬件设计}{楼里信标与展示样机}
\lead{\textbf{可以做硬件 Demo}\quad 外观模型负责解释产品形态，现成蓝牙信标负责验证广播与手机扫描。两者分别展示，不将外壳或扫描成功称为已完成的定位产品。}
\subsection*{外观概念}
建议采用直径约 70 mm、厚约 20 mm 的白色圆角圆盘，带小面积绿色品牌标识、设备编号和维护二维码。尺寸是设计草案，尚未确定结构与生产工艺。背部预留可拆卸底座，面向墙面或吊顶安装；正式安装方式须经物业确认。
\begin{center}
\begin{tikzpicture}[font=\small]
\filldraw[fill=pale,draw=green,line width=1pt] (0,0) circle (1.25);
\draw[line] (0,0) circle (1.08);
\node[green,font=\large\bfseries] at (0,.25) {楼里};
\node[muted,font=\scriptsize] at (0,-.35) {L01\quad 概念外观};
\draw[-{Latex},muted] (1.35,.6)--(2.4,.6);
\node[anchor=west,align=left,muted] at (2.5,.6) {白色外壳与可拆底座\\编号对应部署台账};
\draw[-{Latex},muted] (1.15,-.7)--(2.4,-.7);
\node[anchor=west,align=left,muted] at (2.5,-.7) {维护二维码\\用于查设备与更换记录};
\end{tikzpicture}
\end{center}
\subsection*{功能设计与职责}
信标广播设备标识，手机结合多个观测和场馆地图估计位置；信标本身不计算患者坐标。候选广播标识包括 UUID、Major、Minor，编号与楼层、安装点关联。发射功率、广播间隔、电池方案须随选型和场地验证确定，暂不标注续航年限与定位精度。
\tbl{p{25mm}YY}
\textbf{版本}&\textbf{做什么}&\textbf{能证明什么}\\\hline
外观展示样机&外壳、标识、底座或桌面模型&产品形态与部署设想\\
广播验证样机&现成信标与配置工具，扫描界面&特定手机和场地能接收到广播\\
现场导航样机&多点部署、地图、手机端软件&既定条件下的导航任务表现\\
\end{tabularx}\endgroup\par
\subsection*{小范围布设与维护}
在教学楼走廊或经授权的展示场地，先配置少量设备，放在入口、分岔处与目的地附近。数量和间距按现场测试调整。台账记录编号、位置、配置、安装日期和电池更换；没有网关或在线回传时，设备状态依靠巡检或现场扫描确认。
\subsection*{供应链路线}
首轮向现成信标供应商询问样品、配置方式、协议兼容和供货条件；外观定制与 OEM 留到需求明确后。开发板也可用于固件验证，Nordic 的 nRF52 DK 是一种候选开发平台。[3] 本稿未选定供应商，未发生采购。

\page{03\quad 软件设计}{患者端与场馆管理端}
\lead{\textbf{展示核心}\quad 患者在每个节点都能回答：下一步做什么、去哪里、如何过去、到达后做什么。}
\subsection*{患者端的六个关键画面}
\tbl{p{29mm}Y}
\textbf{画面}&\textbf{内容与交互}\\\hline
开始使用&扫码进入；展示场馆与楼层；可手动选当前位置\\
我的就诊&列出院方确认的任务、状态和必要先后关系；演示使用虚构任务\\
前往下一站&显示目的地、预计步行时间、地标式路线与到达操作\\
途中引导&在关键路口提示方向；偏离时提示重新规划或重新确认位置\\
候诊与返回&显示队列来源和更新时间；用模拟变化解释何时返回\\
完成与反馈&结束当前任务、查看下一站、反馈路线或科室信息错误\\
\end{tabularx}\endgroup\par
\subsection*{体验示例}
一位首次来院的患者完成初诊，需要去检验和影像，再回诊室。演示首先显示任务列表，再打开下一站路线，使用科室和电梯厅作地标；模拟队列变化时展示返回提示，最后呈现复诊和取药。真实检查准备要求及可调整顺序由院方规则决定，导航不提供诊疗判断。
\subsection*{场馆管理端概念}
首版设计包括地图与科室维护、可通行区域、设备部署台账、流程模板和反馈处理。可将临时封闭的通道设为不可通行；无障碍路线按场馆确认的电梯与通道配置。管理端当前未实现，展示时应标注概念界面。
\subsection*{数据与系统接入分阶段}
\begin{enumerate}
\item \textbf{概念演示}\quad 使用虚构地图、任务、队列和位置，不读取真实个人数据。
\item \textbf{路线试点}\quad 使用授权地图、人工配置流程与固定起点或扫码点，观察找路体验。
\item \textbf{定位与业务试点}\quad 验证信标和手机兼容；取得授权后接入必要队列数据，再评估个性化流程。
\end{enumerate}
\subsection*{异常状态也属于产品}
信号不足时允许重新选起点；没有实时队列时显示来源与更新时间；地图路径不可用时提示联系现场导诊；后台数据不可用时回退到静态路线。演示应包括一次回退，解释实际使用的连续性。

\page{04\quad 展示设计}{两分钟讲清产品价值}
\lead{\textbf{这轮先确定展示脚本}\quad 复用已有网页仿真承载医院故事，另做硬件概念说明。下轮按脚本改造展示界面，避免把技术控制台作为默认开场。}
\subsection*{建议的展示结构}
主屏显示患者任务与地图。侧边说明“当前步骤”和“为什么先去这里”。主持人用隐藏的演示控制台切换位置、队列和路线异常；商业观众默认看到服务体验，技术合作方可另行打开算法对比。
\tbl{p{22mm}YY}
\textbf{时间}&\textbf{现场操作}&\textbf{一句解释}\\\hline
0 至 20 秒&打开医院场景，显示虚构就诊任务&第一次来院，患者要把多个地点串起来\\
20 至 55 秒&进入下一站导航，展示一个路口提示&用熟悉的地标告诉他怎样走\\
55 至 85 秒&触发模拟队列变化与返回提醒&把走路和等待放在同一个流程里\\
85 至 105 秒&展示信标外观和地图上的部署点&这些参考点帮助手机确认室内位置\\
105 至 120 秒&显示试点范围与合作请求&从一个楼层验证，再讨论正式部署\\
\end{tabularx}\endgroup\par
\subsection*{演示真实性的显示方式}
网页顶部保留“仿真演示”标识；队列标“模拟队列”；硬件图标“概念外观”。如展示真实广播扫描，单独标注设备和手机条件。完整导航仍为仿真时，不能将扫描画面解释成真机导航已实现。
\subsection*{现有仓库的使用方式}
\begin{itemize}
\item 已有医院、商场、导航和模拟排队，可用于解释基本体验。[1]
\item 当前算法误差是仿真统计，放入技术附录或按需展示。
\item 下一轮拟增加：医院默认开场、预设故事、一键演示、硬件部署说明与试点合作收尾画面。这些属于待开发设计。
\end{itemize}
\subsection*{现场演示的验收条件}
无需讲解算法即可理解流程；故事可重置并重复演示；屏幕能区分模拟和真实来源；两分钟内有明确合作请求。备用方案是按同一脚本播放录屏或展示截图，不依赖现场网络和蓝牙一定可用。
\subsection*{讲给不同观众听}
院方讲流程和服务；集成商讲模块和边界；硬件伙伴讲配置与维护；IDP 评委讲客户发现、付费验证与渠道。所有版本共享同一产品状态说明。

\page{05\quad 商业方案}{从试点到部署与维护}
\lead{\textbf{先验证付费路径}\quad 用小范围项目获得真实需求、交付成本和决策流程，随后决定报价与商业模式。}
\subsection*{拟议服务组合}
\tbl{p{27mm}YY}
\textbf{阶段}&\textbf{客户得到什么}&\textbf{收费与边界}\\\hline
场景调研&访谈、流程梳理、地图勘察与试点范围&可免费访谈；深入勘察另行确定费用\\
单层试点&约定路线演示、测试与评估报告&优先谈范围明确的付费试点；免费试点需设投入上限\\
正式部署&建图、配置、测试、培训与验收&按楼层、接入范围和现场工作量报价\\
年度维护&地图更新、设备巡检或更换、软件支持&明确频次、响应时间与包含的工作量\\
渠道模块&向集成商提供流程导航能力&协商模块费、项目分成和售后责任\\
\end{tabularx}\endgroup\par
\subsection*{首轮报价如何构成}
报价拆为设备与安装、地图与路线配置、软件与业务规则、系统接入、验证培训、维护支持。勘察和接口确认前，只给工作范围和报价结构。现成设备价格、团队工时和合作方分工都需要真实询价，不能把仓库中的假设单价当成交付报价。
\subsection*{每个项目都应记录的经济账}
项目毛贡献按“合同收入减设备、外包、差旅与现场交付直接成本”计算；同时记录内部人力投入，判断报价能否覆盖持续工作。回款节点与现金支出也要单列，防止利润假设掩盖现金压力。
\subsection*{初期获客路线}
通过校友、导师或既有合作方争取业务访谈；由业务负责人指出具体问题，再与信息科确认实施边界。同步访谈医疗信息化集成商，看是否能进入其项目方案。首次会谈以获得对接人和下一次现场讨论为目标，不把口头兴趣计入订单。
\subsection*{可形成的优势与待验证假设}
我们拟围绕院方流程配置、路线维护与实施协作形成能力。优势能否成立，取决于实际交付质量、客户反馈和合作资源。目前不宣称算法独家、硬件专利或已有客户规模。
\subsection*{商业证据逐步累积}
业务访谈记录 → 明确试点意向 → 范围与责任约定 → 验证报告 → 付费与后续部署。每一步都应能提供对应证据，而不是用页面完成度替代客户承诺。

\page{06\quad 试点提案}{一个楼层的合作如何开展}
\lead{\textbf{拟议范围}\quad 一个授权楼层、一段经院方确认的流程、两至四个目的地。试点周期建议为范围确认后的两周，具体排期按场地与参与者条件确定。}
\subsection*{双方投入}
\tbl{p{23mm}Y}
\textbf{楼里}&流程和地图配置、演示原型、测试组织、反馈整理与评估报告；如包含定位验证，另行明确设备与手机端准备\\\hline
\textbf{场地方}&指定业务对接人，提供获授权的地图与科室信息，确认可走路线和安装条件，协调测试时段与参与者\\\hline
\textbf{双方确认}&是否收费、成本上限、测试范围、数据处理、设备拆除、成果使用，以及试点结束后的决策方式\\
\end{tabularx}\endgroup\par
\subsection*{路线验证与定位验证分别评估}
首轮可先测试固定起点的路线与流程体验。若硬件与手机端准备不足，先完成这一部分；信标定位测试作为后续独立阶段。导航体验有效，并不能自动证明定位精度达标。
\subsection*{拟议评估方法}
先用现有标识和指路方式完成一组任务，再用原型完成难度相近的任务；尽量让两种条件的顺序交替，以减少熟悉路线的影响。可先邀请 10 至 20 名志愿者做探索性测试，记录任务完成时间、求助次数、错误转向及主观反馈。样本只能用于发现问题，不外推为全院节省比例。
\subsection*{定位阶段额外记录什么}
记录手机型号、持机方式、场地与人流条件、设备配置、已知参考点和失败情况。分别评估楼层或区域识别、路口引导是否正确，以及适用条件下的误差分布。没有真机记录前，不承诺“全场 1 至 3 米”。
\subsection*{两周的建议安排}
第 1 至 3 天确认范围与路线；第 4 至 6 天配置并做内部预演；第 7 至 10 天组织任务测试；第 11 至 14 天复盘并提交继续、调整或停止的建议。生产系统接入不纳入这一默认排期。
\subsection*{成功与下一步}
双方在开始前约定关心的指标、可接受问题和继续条件。试点完成的交付物包括范围与配置记录、测试结果、问题清单及下一阶段估算；继续部署、补充验证或结束试点都应有明确依据。

\page{07\quad 洽谈工具}{开场介绍与常见问题}
\subsection*{三十秒开场}
“我们在做楼里医院就诊导航，希望把找科室、去检查和等叫号串成一条清楚的路线。现在有可展示的网页仿真，也有信标和软件的产品设计。我们想先了解贵院一个楼层里最困扰患者的找路流程，再讨论一个范围可控的试点，看它是否值得继续投入。”
\subsection*{对方问产品做到哪里了}
“目前能展示流程、路线和模拟队列的交互。真实室内定位、正式队列接入和管理端还要分阶段验证。我们可以先确认服务场景和试点范围，再一起确定优先做哪部分。”
\subsection*{对方问硬件是不是自研}
“楼里信标是我们的产品概念。首轮计划用现成信标验证广播和部署，再决定外观定制、OEM 或进一步研发。目前没有宣称已完成自主芯片、量产或硬件认证。”
\subsection*{对方问精度、价格和交付时间}
“网页精度来自仿真；真实表现需要在指定手机和场地测试。价格与时间取决于楼层、流程和接口范围。我们希望先做一次勘察，形成清楚的试点范围和报价，再讨论交付。”
\subsection*{对方说已有导航系统}
“那我们先看现有系统在哪些环节还需要帮助。如果基础找路已经解决，可以重点讨论就诊任务串联、路线维护或流程体验；如果没有明确缺口，就不建议重复建设。”
\subsection*{对方愿意合作时的收尾}
“能否请您推荐一位门诊业务负责人和信息科对接人，安排一次现场讨论？我们会带着这份方案，先把楼层、流程和评估方式定下来，再提交具体试点范围。”
\subsection*{IDP 汇报应该突出什么}
先说明患者与付费方的不同需求，再展示产品体验、合作方式和商业验证路径。汇报当前学习到的事实、还没验证的假设，以及下一轮将怎样取得证据。以访谈和试点承诺说明推进程度，不把仿真当实测、不把兴趣当销售收入。
\subsection*{首轮会后记录}
记录对象与角色、现有解决方法、最关心的问题、决策链、预算来源、可提供的资源、下一步行动与负责人。合作方尚未确认的部分单独记录，避免形成未经同意的承诺。

\page{08\quad 产品状态与行动}{先完成销售材料再逐步验证}
\subsection*{当前状态与设计范围}
\tbl{p{35mm}Y}
\textbf{已有网页仿真}&商场与医院、定位对比、路径规划、模拟排队、多点任务；基于仓库源码阅读，尚未在本轮做运行验收\\\hline
\textbf{本稿已完成设计}&产品定位、硬件形态、患者与管理端功能、演示脚本、商业方案、试点范围和洽谈话术\\\hline
\textbf{下一轮展示开发}&医院开场、一键故事演示、硬件部署说明、合作收尾画面与录屏\\\hline
\textbf{待现场验证}&手机兼容、真实定位、路线任务表现、安装维护、供应商与成本\\\hline
\textbf{待业务合作}&院方流程确认、系统授权与接口、付费路径、正式验收与维护责任\\
\end{tabularx}\endgroup\par
\subsection*{建议的推进顺序}
\begin{enumerate}
\item 队友先讨论定位与分工，再从本稿整理对外合作摘要，用于约访和解释方案。
\item 按两分钟脚本改造现有 Demo，形成可重复展示的患者故事。
\item 准备信标外观模型；确认供应商与手机端路线后再安排真实广播样机。
\item 访谈首批业务方与集成商，调整首个试点范围；取得授权后开展现场验证。
\end{enumerate}
\subsection*{首批需要取得的证据}
建议先完成 3 至 5 次有效业务访谈，至少弄清一个真实流程的痛点、现有方案、对接人和下一步条件。会谈数是行动建议；它不代表订单目标已达成。试点意向、经确认的场地范围和付费条件应分别记录。
\subsection*{资料依据}
{\footnotesize
[1] 楼里项目仓库，README 与 demo/louli.html，2026 年 10 月 7 日读取。用于确认现有仿真功能与源文件结构。\par
\url{https://github.com/chinayuren2022-2025/louli-indoor-nav}\par
[2] 浙江大学医学院附属邵逸夫医院，智能导航系统项目招标公告，2019 年 5 月 17 日。用于说明历史采购需求，不用于推断当前预算或合作关系。\par
\url{https://www.srrsh.com/articleInfo/3682}\par
[3] Nordic Semiconductor，nRF52 DK Hardware 官方说明，2026 年 10 月 7 日查阅。用于说明蓝牙开发平台候选，不构成供应商选定。\par
\url{https://docs.nordicsemi.com/r/bundle/ug_nrf52832_dk/page/ug/dk/intro.html}\par}
\vfill
\lead{\textbf{本轮成果的用途}\quad 带着清楚的产品概念去展示，带着可执行的试点提案去谈合作，再用获得的反馈决定下一步开发。}
\end{document}
